% 1-16 题 图 1-9：椭圆，右焦点 F 为极点，极轴沿半长轴方向
% 半长轴 A=3（水平）、半短轴 B=2、焦距 c=sqrt(5)≈2.236（示意比例）
\begin{tikzpicture}[scale=0.95, font=\small]
  \coordinate (O) at (0,0);
  \coordinate (Fp) at (2.236,0);       % 极点/右焦点
  \coordinate (Fm) at (-2.236,0);      % 左焦点
  \coordinate (X) at (3.9,0);          % 极轴参考点（F 右侧）
  \coordinate (P) at (2.7217,0.8412);  % 椭圆上一点（θ=60°）
  % 椭圆
  \draw[thick] (O) ellipse (3 and 2);
  % 长、短轴
  \draw[gray!70] (Fm) -- (Fp);
  \draw[gray!70] (0,-2) -- (0,2);
  % 极轴（从 F 向右）
  \draw[->, >=Stealth, thick, blue!60!black] (Fp) -- (4.15,0) node[right] {极轴};
  % 位矢 r = FP
  \draw[dashed] (Fp) -- (P);
  % θ 角（自极轴量起）
  \pic[draw, angle radius=8mm, angle eccentricity=1.4, "$\theta$"] {angle = X--Fp--P};
  % 标注
  \fill (O) circle (1.8pt) node[below left] {$O$};
  \fill (Fp) circle (2pt) node[below] {$F$};
  \fill (Fm) circle (1.5pt) node[below] {$F'$};
  \fill (P) circle (2pt) node[above right] {$P$};
  \node[below] at (1.5,0) {$A$};        % 半长轴
  \node[right] at (0,1.05) {$B$};       % 半短轴
\end{tikzpicture}
